\section{Zadanie 7}
\autor{Hubert Jastrzębski}
Łamiemy kij w dwóch niezależnie i jednostajnie wybranych punktach. Jakie
jest prawdopodobieństwo tego, że z otrzymanych kawałków można zbudować trójkąt?


\subsection*{Rozwiązanie}
Niech kij ma długość \(1\) i zostanie złamany w punktach \(X_1\) i \(X_2\), gdzie \(X_1, X_2 \sim \text{Unif}(0, 1)\).

Funkcje gęstości zmiennych losowych \(X_1\) i \(X_2\) są następujące:

\[
	f_1(x_1) = \frac{1}{1 - 0} = 1 \quad \text{dla } x_1 \in [0, 1] \quad \land \quad f_2(x_2) = \frac{1}{1 - 0} = 1 \quad \text{dla } x_2 \in [0, 1]
\]

Zmienne losowe są niezależne, więc ich funkcja gęstości wspólnej jest iloczynem funkcji gęstości poszczególnych zmiennych:
\[
	f(x_1, x_2) = f_1(x_1) \cdot f_2(x_2) = 1 \cdot 1 = 1 \quad \text{dla } x_1, x_2 \in [0, 1]
\]

Policzmy zdarzenie odwrotne, czyli prawdopodobieństwo, że z kawałków nie można zbudować trójkąta.
Oznacza to, że jeden z kawałków jest dłuższy niż \(0.5\).

Mamy trzy przypadki:

\medbreak
\textbf{1. Pierwszy kawałek jest długi:}

Przez pierwszy kawałek mam na myśli kawałek najbardziej po lewej stronie, który ma długość \(min(X_1, X_2)\).
Aby był długi (\(\geq 0.5\)), musi zachodzić \(X_1 \geq 0.5 \land X_2 \geq 0.5\).

Prawdopodobieństwo na to wynosi:
\[
	P(X_1 \geq 0.5 \land X_2 \geq 0.5) = P(X_1 \geq 0.5) \cdot P(X_2 \geq 0.5)
	= \int_{0.5}^{1} f(x_1) \, dx_1 \cdot \int_{0.5}^{1} f(x_2) \, dx_2 = \frac{1}{2} \cdot \frac{1}{2} = \frac{1}{4}
\]

\medbreak
\textbf{2. Trzeci kawałek jest długi:}

Analogicznie \(P(X_1 \leq 0.5 \land X_2 \leq 0.5) = \frac{1}{4}\)

\medbreak
\textbf{3. Drugi kawałek jest długi:}
\[
	P(|X_1 - X_2| \geq 0.5) = P(X_1 - X_2 \geq 0.5) + P(X_2 - X_1 \geq 0.5)
\]
\[
	= \int \int_{\{x_1 \geq x_2 + 0.5\}} f(x_1, x_2) \, dx_1 \, dx_2 + \int \int_{\{x_2 \geq x_1 + 0.5\}} f(x_1, x_2) \, dx_1 \, dx_2
\]

Rozważmy \(\int \int_{\{x_1 \geq x_2 + 0.5\}} f(x_1, x_2) \, dx_1 \, dx_2\):

Pamiętajmy, że \(x_1, x_2 \in [0, 1]\), więc \(x_2 + 0.5 \leq 1 \implies x_2 \leq 0.5\).

Natomiast dla ustalonego \(x_2\), \(x_1\) może przyjmować wartości od \(x_2 + 0.5\) do \(1\). Tak więc:

\[
	\int \int_{\{x_1 \geq x_2 + 0.5\}} f(x_1, x_2) \, dx_1 \, dx_2 = \int_{0}^{0.5} \int_{x_2 + 0.5}^{1} f(x_1, x_2) \, dx_1 \, dx_2
\]
\[
	= \int_{0}^{0.5} \int_{x_2 + 0.5}^{1} 1 \, dx_1 \, dx_2 = \int_{0}^{0.5} (1 - (x_2 + 0.5)) \, dx_2 = \int_{0}^{0.5} (0.5 - x_2) \, dx_2
\]
\[
	= \left[0.5x_2 - \frac{x_2^2}{2}\right]_{0}^{0.5} = \left(0.5 \cdot 0.5 - \frac{(0.5)^2}{2}\right) - (0) = \frac{1}{4} - \frac{1}{8} = \frac{1}{8}
\]

Analogicznie \(\int \int_{\{x_2 \geq x_1 + 0.5\}} f(x_1, x_2) \, dx_1 \, dx_2 = \frac{1}{8}\).

Tak więc \(P(|X_1 - X_2| \geq 0.5) = \frac{1}{8} + \frac{1}{8} = \frac{1}{4}\).

\medbreak
Sumując wszystkie przypadki mamy:
\[
	P(\text{jeden kawałek jest długi}) = \frac{1}{4} + \frac{1}{4} + \frac{1}{4} = \frac{3}{4}
\]
\[
	P(\text{żaden kawałek nie jest długi}) = 1 - \frac{3}{4} = \frac{1}{4}
\]

\subsection*{Odpowiedź}
\[
	\boxed{\frac{1}{4}}
\]
